\documentclass[12pt]{article}

%% Basic document formatting
\usepackage{amsmath, amsthm, amssymb, amsfonts}
\usepackage{mathtools}
\usepackage{xspace}
\usepackage{thmtools}
\usepackage{graphicx}
\usepackage{setspace}
\usepackage{fancyhdr}
\usepackage{titling}
\usepackage{hyperref}
\usepackage[left=0.4in,right=0.4in,top=1in,bottom=1in]{geometry}
\usepackage{float}
\usepackage{tabularx}
\usepackage[utf8]{inputenc}
\usepackage[english]{babel}
\usepackage{framed}
\usepackage[dvipsnames]{xcolor}
\usepackage{environ}
\usepackage{tcolorbox}
\tcbuselibrary{theorems,skins,breakable}

\usepackage{enumitem}
\setlist[enumerate]{leftmargin=*}
\setlist[enumerate,1]{labelindent=\parindent}
\setlist[enumerate,2]{labelindent=0pt}

% Blackboard Bold
\newcommand{\Z}{\mathbb{Z}}                 % integers
\newcommand{\Zpl}{\mathbb{Z}_{+}}           % positive integers
\newcommand{\Znn}{\mathbb{Z}_{\geq 0}}      % nonnegative integers. 
\newcommand{\Q}{\mathbb{Q}}                 % rationals
\newcommand{\Qpl}{\mathbb{Q}_{+}}           % positive Rationals
\newcommand{\R}{\mathbb{R}}                 % reals
\newcommand{\Rpl}{\mathbb{R}_{+}}           % positive Reals
\newcommand{\C}{\mathbb{C}}                 % complex numbers
\newcommand{\F}{\mathbb{F}}                 % field

% Words
\newcommand{\st}{\text{ such that }}
\newcommand{\wrt}{\text{ with respect to }}
\newcommand{\with}{\text{ with }}
\newcommand{\ie}{\text{, i.e. }}
\newcommand*{\ditto}{\text{---}\texttt{"}\text{---}}

% Operators
\newcommand{\abs}[1]{\left|#1\right|}                   % absolute value
\newcommand{\norm}[1]{\abs{\abs{#1}}}
\newcommand{\floor}[1]{\left\lfloor #1 \right\rfloor}   % floor
\newcommand{\ceil}[1]{\left\lceil #1 \right\rceil}      % ceiling
\newcommand{\powset}[1]{\mathscr{P}(#1)}

% Category Theory 
\renewcommand{\hom}{\text{Hom}}
\newcommand{\homspace}[3]{\hom_{#1}(#2, #3)}
\newcommand{\coproduct}[9]{
  \begin{tikzcd}[ampersand replacement=\&]
      \& \& #1 \arrow[dll, bend right, "#6"] \arrow[dl, "#5"] \\
   #4 \& #3 \arrow[l, "#9"] \\
      \& \& #2 \arrow[ull, bend left, "#8"] \arrow[ul, "#7"]
  \end{tikzcd}
}

\newcommand{\product}[9]{ 
  \begin{tikzcd}[ampersand replacement=\&]
      \& \& #3 \\
      #1 \arrow[r, "#7"] \arrow[urr, bend left, "#5"] \arrow[drr, bend right, "#6"] \& #2 \arrow[ur, "#8"] \arrow[dr, "#9"] \\ 
      \& \& #4
  \end{tikzcd}
}


% Algebra 
\newcommand{\Syl}[2]{\text{Syl}_{#1}{(#2)}}           % set of Sylow-p subgroups 
\newcommand{\<}{\langle}                % \<x\>, subgroup generated by x 
\renewcommand{\>}{\rangle}
\renewcommand{\index}[2]{[#1 : #2]}
\newcommand{\id}{\text{id}}             % identity element
\newcommand{\lhdneq}{%
  \mathrel{\ooalign{$\lneq$\cr\raise.22ex\hbox{$\lhd$}\cr}}}
\newcommand{\order}[1]{\text{o}(#1)}    % order of an element
\newcommand{\spec}{\operatorname{Spec}}
\newcommand{\rad}{\operatorname{rad}}   % radical of an ideal 
\newcommand{\sym}[1]{\mathfrak{S}_{#1}}
\newcommand{\aut}[1]{\text{Aut}(#1)} 
\newcommand{\gal}[2]{\text{Gal}(#1 / #2)} 
\newcommand{\inn}[1]{\text{Inn}(#1)} 
\let\oldcong\cong
\let\oldequiv\equiv
\renewcommand{\cong}{\oldequiv}
\renewcommand{\equiv}{\oldcong}

\newcommand{\triangleleftneq}{\mathrel{\ooalign{%
  $\trianglelefteq$\cr
  \hidewidth\raise0.06ex\hbox{$\not\mkern4mu$}\hidewidth\cr
}}}

\makeatletter % cycle
\newcommand{\cyc}[1]{(\mathbf{\cyc@process#1\relax})}
\def\cyc@process#1#2\relax{%
	#1%
	\ifx\relax#2\relax
	\else
		\,\cyc@process#2\relax
	\fi
}
\makeatother

\newcommand{\GL}[2]{\text{GL}_{#1}(#2)}                             % general linear group
\newcommand{\SL}[2]{\text{SL}_{#1}(#2)}                             % special linear group
\newcommand{\gl}[2]{\mathfrak{gl}_{#1}(#2)}
\renewcommand{\sl}[2]{\mathfrak{sl}_{#1}(#2)}
\newcommand{\so}[2]{\mathfrak{so}_{#1}(#2)}
\newcommand{\su}[1]{\mathfrak{su}(#1)}

% Linear Algebra
\newcommand{\bmat}[1]{\begin{bmatrix}#1\end{bmatrix}}   % bracketed matrix
\newcommand{\rank}{\operatorname{rank}}                 % rank
\newcommand{\nullity}{\operatorname{nullity}}           % nullity
\newcommand{\tr}{\operatorname{tr}}

% Combinatorics
\newcommand{\qnum}[1]{\left[#1\right]_q}                % q-analog
\newcommand{\qfact}[1]{\left[#1\right]!_q}              % q-analog factorial 
\newcommand{\qbinom}[2]{\binom{#1}{#2}_q}               % Gaussian binomial coefficient

% Topology/Analysis
\newcommand{\ball}[2]{\text{B}_{#1}(#2)}  % B_r(x): open r-balls around x
\newcommand{\diam}{\text{diam}}           % diamter of a set in metric space 
\newcommand{\lowint}[2]{\underline{\int_{#1}^{#2}}}
\newcommand{\upint}[2]{\overline{\int}_{#1}^{#2}}

% Misc Notation
\newcommand{\oneton}[1]{\{1, \ldots, #1\}}
\newcommand{\defeq}{\vcentcolon=}   % :=
\newcommand{\eqdef}{=\vcentcolon}   % =: 
\renewcommand{\bf}[1]{\textbf{#1}}
\newcommand{\im}{\operatorname{im}}
\newcommand{\fnrest}[2]{\left.#1\right|_{#2}}
\newcommand{\isom}{\overset{\sim}{\rightarrow}} 


\renewcommand{\thesection}{\arabic{section}.}
\renewcommand{\thesubsection}{(\alph{subsection})}

\newtheorem{claim}{Claim}
\newtheorem*{lemma}{Lemma}

%% Headers & title setup
\newcommand{\course}{Math100B}
\newcommand{\myname}{Jay Ser}
\setlength{\headheight}{14.5pt}
\pagestyle{fancy}
\fancyhf{}
\renewcommand{\headrulewidth}{0.4pt}
\lhead{\course}
\rhead{\myname}
\cfoot{\thepage}
\setlength{\droptitle}{-4em} 
\title{\course\ - HW \#7}
\author{\myname}
\date{2026.03.01}

\begin{document}
\maketitle
\thispagestyle{fancy}

%------------------------------------------------------------------------------%

\section{} 
\subsection{}
I leave the matrix computations mod $p$ to Gavin. 
To solve the systems, I row reduce and represent row redution via elementary matrices. 
When $p = 5$, 
\begin{align*}
  \bmat{1 & 0 \\ 0 & 3} \bmat {1 & -1 \\ 0 & 1} \bmat{1 & 0 \\ -2 & 1} \bmat{6 & -3 & 3 \\ 2 & 6 & 1} & = \bmat{1 & 0 \\ 0 & 3} \bmat {1 & -1 \\ 0 & 1} \bmat{1 & 0 \\ -2 & 1} \bmat{1 & 2 & 3 \\ 2 & 1 & 1} \\ 
          & = \bmat{1 & 0 \\ 0 & 3} \bmat {1 & -1 \\ 0 & 1} \bmat{1 & 2 & 3 \\ 0 & 2 & 0} \\ 
          & = \bmat{1 & 0 \\ 0 & 3} \bmat{1 & 0 & 3 \\ 0 & 2 & 0} \\ 
          & = \bmat{1 & 0 & 3 \\ 0 & 1 & 0}
\end{align*}
So $X = \bmat{3 \\ 0}$ is the unique solution to the system $AX = B$. 

When $p = 7$, 
\begin{align*}
  \bmat{1 & 0 \\ -2 & 1} \bmat {6 & 0 \\ 0 & 1} \bmat{6 & -3 & 3 \\ 2 & 6 & 1} & = \bmat{1 & 0 \\ -2 & 1} \bmat{1 & 3 & 4 \\ 2 & 6 & 1} \\ 
          & = \bmat{1 & 3 & 4 \\ 0 & 0 & 0}
\end{align*}
So any $X = \bmat{x_1 \\ x_2}$ satisfying $x_1 + 3x_2 = 4$ is a solution to the system $AX = B$. 

When $p = 11$, 
\begin{align*}
  \bmat{1 & -3 \\ 0 & 1} \bmat{1 & 0 \\ -6 & 1} \bmat{6 & 0 \\ 0 & 1} \bmat{0 & 1 \\ 1 & 0} \bmat{6 & -3 & 3 \\ 2 & 6 & 1} & = \bmat{1 & -3 \\ 0 & 1} \bmat{1 & 0 \\ -6 & 1} \bmat{6 & 0 \\ 0 & 1} \bmat{2 & 6 & 1 \\ 6 & 8 & 3} \\ 
          & = \bmat{1 & -3 \\ 0 & 1} \bmat{1 & 0 \\ -6 & 1} \bmat{1 & 3 & 6 \\ 6 & 8 & 3} \\ 
          & = \bmat{1 & -3 \\ 0 & 1} \bmat{1 & 3 & 6 \\ 0 & 1 & 0} \\ 
          & = \bmat{1 & 0 & 6 \\ 0 & 1 & 0} 
\end{align*}
So $X = \bmat{6 \\ 0}$ is the unique solution to the system $AX = B$. 

\subsection{}
Using the formula for the determinant of a $3 \times 3$ matrix, 
$$\det{\bmat{1 & 2 & 0 \\ 0 & 3 & -1 \\ -2 & 0 & 2}} = 10$$ 
Since reduction mod $p$ respects addition and multiplication, the determinant of the matrix mod $p$ is $10 \pmod{p}$. 
Hence the matrix is invertible mod $p \iff p \neq 2, 5$. 

\subsection{} 
Over any of the fields $\Q, \F_2, \F_3$, the following row reduction makes sense: 
\begin{equation} \label{eq:res_mat}
  \bmat{1 & -1 & \\ & 1 & \\ & & 1} \bmat{1 & & & \\ -1 & 1 &  \\ & & 1} \bmat{1 & & & \\  & 1 &  \\ -1 & & 1} \bmat{1 & 1 & 0 & 1 \\ 1 & 0 & 1 & -1 \\ 1 & -1 & -1 & 1} = \bmat{1 & 0 & 1 & -1 \\ 0 & 1 & -1 & 2 \\ 0 & 2 & 1 & 0}  
\end{equation}
Over $\F_{2}$, the resulting matrix in \eqref{eq:res_mat} is 
$$\bmat{1 & & 1 & 1 \\ & 1 & 1 & \\ & & 1 }$$
It can be row-reduced to 
$$\bmat{1 & &  \\  & 1 & -1 \\  &  & 1} \bmat{1 & & -1\\  & 1 &  \\  &  & 1} \bmat{1 & & 1 & 1 \\ & 1 & 1 & \\ & & 1 } = \bmat{1 & & & 1 \\ & 1 & & 0 \\ & & 1 & 0}$$
Hence $X = \bmat{1 \\ 0 \\ 0}$ is the unique solution to $AX = B$ over $\F_{2}$. 

Over $\F_3$, the resulting matrix in \eqref{eq:res_mat} is 
$$\bmat{1 & & 1 & 2 \\ & 1 & 2 & 2 \\  & 2 & 1 & }$$
It can be row-reduced to 
$$\bmat{1 & & \\ & 1 & \\ & -2 & 1} \bmat{1 & & 1 & 2 \\ & 1 & 2 & 2 \\  & 2 & 1 & } = \bmat{1 & & 1 & 2 \\ & 1 & 2 & 2 \\ & & & 1}$$
The bottom row indicates that the system $AX = B$ has no solutions over $\F_3$. 

Over $\Q$, row-reduce \eqref{eq:res_mat} as such:  
$$\bmat{1 & & -1 \\ & 1 & \\ & & 1} \bmat{1 & & \\ & 1 & 1 \\ &  & 1}\bmat{1 & & \\ & 1 & \\ &  & \frac{1}{3}} \bmat{1 & & \\ & 1 & \\ & -2 & 1} \bmat{1 & 0 & 1 & -1 \\ 0 & 1 & -1 & 2 \\ 0 & 2 & 1 & 0} = \bmat{1 & & & \frac{1}{3} \\ & 1 & & \frac{2}{3} \\ & & 1 & -\frac{4}{3}}$$
So $X = \bmat{1/3 \\ 2/3 \\ - 4/3}$ is the unique solution to $AX = B$ over $\Q$. 

\subsection{}
Notice that 
$$\bmat{1 & 1 \\ & 1}^n = \bmat{1 & n \\ & 1}$$
The identity in $\GL_2(\F_7)$ is just the $2 \times 2$ identity matrix $I_2$, so the order of $\bmat{1 & 1 \\ & 1}$ in $\GL_2(\F_7)$ is 7. 

\section{} 
\subsection{}
$0v = (0 + 0)v = 0v + 0v \implies 0v = 0$ by the distributive property of scalar multipliation and the cancellation law. 

\subsection{}
$0 \in F$, so if $w \in W$, $0w = 0 \in W$. 

\subsection{} 
$-1(w) + w = (-1 + 1)w = 0w = 0$ by the distributive property of scalar addition, so $-1(w)$ is the additive inverse of $w$. 
Hence $-w = -1(w) \in W$ if $w \in W$. 

\section{}
\subsection{}
It is clear that $P_d$ is an abelian group under the usual polynomial addition since $F$ is a field and the sum of polynomials of degree at most $d$ has degree at most $d$ as well. 
Let scalar multiplication in $P_d$ be the usual polynomial multiplication, where the scalar is treated as a degree 0 polynomial. 
Take any $f(x), g(x) \in F[x]$ with degree at most $d$.
By properties of the ring of polynomials, $\forall \lambda, \mu \in F \setminus 0$,  
$$\lambda(f(x) + g(x)) = \lambda f(x) + \lambda g(x) \text{ and } (\lambda + \mu)f(x) = \lambda f(x) + \mu f(x)$$ 
where clearly the degrees of $\lambda f(x) + \lambda g(x)$ and $\lambda f(x) + \mu f(x)$ are at most $d$. 
This shows that $P_d$ is a vector space over $F$. 

\subsection{}
$\{1, x, \ldots, x^d\}$ form a basis for $P_d$. 
By definition of an indeterminant, 
$$\alpha_0 1 + \alpha_1 x + \ldots + \alpha_d x^d = 0 \iff \alpha_0 = \alpha_1 = \ldots = \alpha_d = 0$$
so $\{1, x, \ldots, x^d\}$ is a linearly independent set.   
It also follows by definition of the polynomial ring that any polynomial in $P_d$ is a linear combination of $\{1, x, \ldots, x^d\}$.
This shows that $\{1, x, \ldots, x^d\}$ is a basis for $P_d$. 

\section{} 
\subsection{}
Because matrix addition is pointwise and $F$ is a field, namely an additive abelian group, matrix multiplicaiton clearly makes $M_n(F)$ into an additive abelian group where the zero matrix, which I will denote $[0]$, is the additive identity.  
Also, define the mutiplicaiton of a scalar $\lambda \in F$ with a matrix to be pointwise multiplication;
that is, multiply every entry of the matrix by $\lambda$. 
Once again, because $F$ is a field, pointwise multiplication by a scalar and pointwise addition satisfies the two distributive properties for a vector space. 

\subsection{}
Because matrix multiplicaiton and scalar addition is pointwise, the sum of two symmetric matrices is a symmetric matrix and a symmetric matrix multiplied by a scalar is a symmetric matrix, 
i.e., the set of all symmetric matrices is closed under addition and scalar multiplication, hence it is a subspace of $M_n(F)$. 

\subsection{}
$\{s_{ij}\}_{1 \leq i \leq j \leq n}$, where $s_{ij}$ is a $n \times n$ matrix with $1$ in the $i, j$th entry and the $j, i$th entry and zeros everywhere else, form a basis for $n \times n$ symmetric matrices.  
By definition, each $s_{ij}$ is symmetric, and it is clear that the $s_{ij}$'s are linearly independent; 
for example, the sum of any linear combination of the $s_{ij}$'s is zero if and only if the coefficient of each term is 0 since each $s_{ij}$ occupies entries that are occupied my no other $s_{i'j'}$.  
It is also clear that the $s_{ij}$'s span the space of all symmetric matrices since any $i, j$th entry of a symmetric matrix can be represented uniquely by a scalar multiple of $s_{ij}$. 

\section{}
\subsection{}
One aleady knows that given any $A \in M_{m \times n}(F)$, any $X, Y \in F^n$, and any $\lambda \in F$,
$$A(X + Y) = AX + AY, \: A(\lambda X) = (\lambda A) X$$
So it follows immediately that $\varphi_{A}$ is a linear transformation. 

\subsection{} 
By above, $\varphi_{A}: F^n \rightarrow F^m$ defined by matrix multiplication is a linear transformation. 

Conversely, given any linear transformation $\varphi: F^n \rightarrow F^m$, let 
$$A = \bmat{\varphi(e_1) & \varphi(e_2) & \ldots & \varphi(e_n)}$$ 
be a $m \times n$ matrix where the $\varphi(e_i)$'s are written as coordinate vectors with respect to the standard basis in $F^m$. 
Recall that $Ae_i$ gives the $i$th column of $A$.
Thus for any $v = \bmat{v_1 \\ \vdots \\ v_n} \in F^n$ written as a coordinate vector with respect to the standard basis in $F^n$, $v = v_1 e_2 + v_2 e_2 + \ldots v_n e_n$ implies  
\begin{align*}
  Av  & = v_1 Ae_1 + v_2 Ae_2 + \ldots + v_n Ae_n \\ 
      & = v_1 \varphi(e_1) + v_2 \varphi(e_2) + \ldots + v_n \varphi(e_n) \\ 
      & = \varphi(v)
\end{align*}
This shows that linear transformations and matrices are different representations of the same thing.  

\subsection*{(c) and (d)}
(c) follows immediately from (d) since the fact that $\varphi^{-1}$ can be represented by a matrix automatically means $\varphi^{-1}$ is a linear transformation. 

(d) is obvious too. 
$\psi$ is the inverse of $\varphi$ if and only if $m = n$ \textit{(quick proof of this in case it cannot be assumed: if $m > n$, then $\varphi$ isn't surjective since the images of the basis elements for $F^n$ form a basis for the $n$-dimensional subspace $\varphi(F^n)$ of $F^m$
Similarly, if $m < n$, then $\varphi$ is not injective)} and  
\begin{equation} \label{eq:fn_inv} 
  \forall v \in F^n, \, \psi(\varphi(v)) = v \text{ and } \forall w \in F^m, \, \varphi(\psi(w)) = w
\end{equation}
If $B$ is the matrix associated to $\psi$ and $A$ is the matrix associated with $\varphi$, then \eqref{eq:fn_inv} is equivalent to 
\begin{equation} \label{eq:mat_inv}
  \forall v \in F^n, \, BAv = v \text{ and } \forall w \in F^m, \, ABw = w
\end{equation}
and, \eqref{eq:mat_inv} is of course equivalent to saying $BA = I_n$ and $AB = I_m$. 
But $m = n$, so $B = A^{-1}$.  

\section{}
This follows from Q5. 
Let $A = \bmat{v_1, \ldots, v_n}$. 
If $\{v_1, \ldots, v_n\}$ is a basis for $F^n$, then any $Y \in F^n$ is represented by a unique linear combination of $v_1, \ldots, v_n$.  
In other words, the linear transformation $\varphi_A$ associated to $A$ is surjective and injective, hence bijective. 
This means that $A$ is invertible. 

Conversely, suppose $\{v_1, \ldots, v_n\}$ is not a basis for $F^n$.
$F^n$ is $n$-dimensional; 
it must be that $\{v_1, \ldots, v_n\}$ do not span $V$ and are not linearly independent (rank nullity theorem or whatever, I think it also follows from what was discussed in lecture? Haven't gone to lecture...).
Namely, if $\{v_1, \ldots, v_n\}$ do not span $V$, then $\exists Y \in F^n \st \forall X \in F^n$, $AX \neq Y$. 
So $A$ is not surjective, which means it is not invertible. 

\section{}
\subsection{}
This follows immediately from Q6. 
Every distinct element of $\GL_2(\F_p)$ corresponds to a distinct (ordered) basis for $\F_{p}^{2}$. 

\subsection{}
Let's count the number of ordered bases for $\F_{p}^{2}$. 
This corresponds to the number of ordered pairs $(v_1, v_2)$ in $\F_{p}^{2}$ where $v_1$ and $v_2$ are linearly independent.
$v_1$ can be chosen to be any element in $\F_{p}^{2}$ except 0;
there are $p^2 - 1$ such elements. 
Since the only requirement for $v_2$ is that $v_2$ is linearly independent from $v_1$, i.e., not a scalar multiple of $v_1$, there are $p^2 - p$ choices for $v_2$. 
In all, this gives 
$$(p^2 - 1)(p^2 - p) = p(p - 1)^2(p + 1)$$
choices for the pair $(v_1, v_2)$, i.e., this is the order of $\GL_{2}(\F_p)$.

\end{document}
